\scanpage{74}
\mspage{25}
\unsure{Let us return|Returning} for a moment to $X$, and
consider the set \deleted{\unsure{of|of the}} $I_0$ of the $i\in I$ such that
$X_i=\varnothing$. This is a sieve, and we have
$X_i^*=\varnothing$, $V_{ij}^*=V_{ij}=\varnothing$ if $i\in I_0$. Thus neglecting
the empty spaces, we see that the diagram
of type $\widetilde I$ defined by the stratified space $X$
does not change essen\textsuperscript{t}ially on replacing $I$ by
$I\setminus I_0$, or more generally by $I\setminus I'_0$, where $I'_0\subset I'$
is a sieve,
which gives rise \unsure{to|:} a diagram
$\widetilde{I\setminus I'_0}$ which is \uline{contained} in $\widetilde I$ (this is
true for \uline{every} sieve of $I$).

If for example we have two sieves
\[
(\unsure{14|16})\quad I''\subset I'\subset I
\]
whence
\[
(15)\quad X_{I''}\subset X_{I'}\quad\text{and}\quad U_{I',I''}\overset{\mathrm{def}}{=}X_{I'}\setminus X_{I''},
\]
we can \deleted{\unsure{on}} (and we \unsure{will|have}) on $U_{I',I''}$ consider
instead the \deleted{filtration} \inserted{stratification} of type $I'\setminus I''$,
defined by the \deleted{$X_{i'}\setminus X_{i'}\cap X_{I''}\quad(i'\in I')\subset X$}
\[
(16)\quad X_{i',\unsure{\deleted{\unsure{X}}\inserted{I',U''}|\unsure{U|X}_{I',U''}}}=X_{i'}\setminus X_{i'}\cap X_{I''}\subset U_{I',I''}
\]
whose elementary topoi are therefore the
$X_{i'}^*$ $(i'\in I'\setminus I'')$ and some $V^*_{i',j'X'}$, $V_{i',j'X'}$

% Thin grey foliation at the lower left, in another hand.
\margin{73}
